\subsection{Définition formelle de la fonction \textsc{Digamma\_Savard}}

\subsubsection{Réduction algébrique : le polynôme $B_0(k)$}

L'observation décisive est que la reconstruction
\[
  P \;=\; \frac{\SB - \mathrm{Digamma}}{k^{6}}
\]
se simplifie : $\SB - \SA$ est divisible par $k^{6}$ pour tout $k$, et le
quotient est un polynôme en $k$ à cinq termes. En effet, au niveau
$n = 10$,
\begin{align*}
  \SA(k,10) &= k^{10}+k^{9}+k^{6}+k^{5}+k^{4}+k^{3}+k^{2}+k,\\
  \SB(k,10) &= k^{11}+k^{10}+k^{7}+k^{5}+k^{4}+k^{3}+k^{2}+k,\\
  \SB(k,10)-\SA(k,10) &= k^{11}-k^{9}+k^{7}-k^{6}
  \;=\; k^{6}\bigl(k^{5}-k^{3}+k-1\bigr).
\end{align*}

\begin{boiteformule}[Polynôme de base $B_0$]
\[
  \boxed{\;B_0(k) \;=\; \frac{\SB(k,10)-\SA(k,10)}{k^{6}}
  \;=\; k^{5}-k^{3}+k-1\;}
\]
\end{boiteformule}

Le Digamma étant $\SA \pm k^{8}$ ou $\SA \pm k^{7}$, on obtient, pour
chacune des quatre branches, un candidat \emph{polynomial} en $k$ :
\begin{align*}
  P_{A(8)-}(k) &= B_0(k) + k^{2}, &
  P_{A(8)+}(k) &= B_0(k) - k^{2},\\
  P_{A(7)+}(k) &= B_0(k) - k, &
  P_{A(7)-}(k) &= B_0(k) + k.
\end{align*}

\begin{boitenote}[Pourquoi la polarité de la suite A commande le signe]
Le signe du Digamma et le signe de la correction \emph{retombent} sur la
polarité de la position dans la suite~A : la position~8 (terme
$k^{8}$) donne une correction en $k^{2}$ \emph{positive} quand elle est
soustraite, la position~7 (terme $k^{7}$) donne une correction en $k$
\emph{positive} quand elle est \emph{ajoutée}. Ce renversement n'est pas
conventionnel : il découle mécaniquement de
$\SB-\SA = k^{6}B_0$, où le terme de rang~6 est celui qui a été
\emph{retiré} à la suite~B par le saut Zêta.
\end{boitenote}

\subsubsection{Table de correspondance : branche, signe et terme}

\begin{center}
\small
\begin{tabular}{@{}l c c r@{}}
\toprule
\rowcolor{spectralgris}
Branche & Position dans la suite A & Terme correctif & Candidat polynomial \\
\midrule
$A(8)-$ & 8\textsuperscript{ième} & $-k^{8}$ & $B_0+k^{2}$ \\
$A(8)+$ & 8\textsuperscript{ième} & $+k^{8}$ & $B_0-k^{2}$ \\
$A(7)+$ & 7\textsuperscript{ième} & $+k^{7}$ & $B_0-k$ \\
$A(7)-$ & 7\textsuperscript{ième} & $-k^{7}$ & $B_0+k$ \\
\bottomrule
\end{tabular}
\end{center}

Cette table est la \emph{vraie} table de correspondance : elle a été
revérifiée numériquement sur $676\,191$ ancres, après qu'une version
antérieure (issue du script de criblage originel) eut inversé les
lignes $A(7)+$ et $A(7)-$. L'inversion était invisible au regard mais
destructrice : elle inversait le signe de la correction Digamma pour
une branche sur quatre.

\subsubsection{Les quatre ancres de référence : $P = 29,\ 31,\ 37,\ 41$}

Pour le rapport typique $t = 1/2$ (soit $k = 2$), les suites à dix termes
valent
\[
  \SA(2,10) = 2+4+8+16+32+64+128+256+384+768 = 1\,662,
\]
\[
  \SB(2,10) = 2+4+8+16+32+128+256+512+768+1536 = 3\,262,
\]
et la 6\textsuperscript{ième} position de la suite A (le « Zêta ») vaut
$2^{6} = 64$.

\begin{center}
\small
\begin{tabular}{@{}c r r r r r r@{}}
\toprule
\rowcolor{spectralgris}
$n$ & $P$ & $\SA(2,n)$ & $\SB(2,n)$ & $\mathrm{Digamma}$ & $\SA+\mathrm{Digamma}$ & Branche \\
\midrule
10 & 29  & 1\,662   & 3\,262   & $-256$   & 1\,406   & $A(8)-$ \\
11 & 31  & 3\,326   & 6\,590   & $+1\,280$ & 4\,606   & $A(8)+$ \\
12 & 37  & 6\,654   & 13\,246  & $+4\,224$ & 10\,878  & $A(8)+$ \\
13 & 41  & 13\,310  & 26\,558  & $+10\,624$ & 23\,934  & $A(8)+$ \\
\bottomrule
\end{tabular}
\end{center}

\begin{boiteformule}[Vérification de la reconstruction sur les quatre ancres]
\[
  \frac{\SB - \mathrm{Digamma}_{\mathrm{calculé}}}{64}
  \;=\;
  \frac{3262-1406}{64}=29,\quad
  \frac{6590-4606}{64}=31,\quad
  \frac{13246-10878}{64}=37,\quad
  \frac{26558-23934}{64}=41.
\]
\end{boiteformule}

\begin{boitenote}[Point d'exactitude relevé à la relecture]
La valeur du Digamma pour $n = 13$ vaut $10\,624$ et donne
$\SA + \mathrm{Digamma} = 23\,934$. La reconstruction
$(26\,558-23\,934)/64 = 41$ est donc exacte. Une saisie antérieure
portait $23\,936$ (soit $10\,626$), qui ne redonne pas $41$ mais
$40{,}96875$ : l'écart est de $2$ sur le Digamma et faute de
reproduction du nombre premier. La valeur $10\,624$ est la bonne.
\end{boitenote}

\subsubsection{La récurrence empirique observée}

Les quatre valeurs du Digamma se lisent aussi comme des décompositions
sur les trois derniers termes de la suite A de référence
($256$, $384$, $768$), avec des coefficients formant une progression
de pas $4$ ($1,\ 5,\ 9,\ 13$) :
\begin{align*}
  256     &= 1\times256,\\
  1\,280  &= 5\times256,\\
  4\,224  &= 9\times256 + 5\times384,\\
  10\,624 &= 13\times256 + 9\times384 + 5\times768.
\end{align*}
Cette décomposition canonique s'écrit
\begin{equation*}
  \mathrm{Digamma}(2,n) \;=\; \sum_{j=0}^{\max(0,\,n-11)}
  \bigl(4\,(n-10-j)+1\bigr)\,b_j,
  \qquad
  b_0 = 256,\; b_1 = 384,\; b_2 = 768,\; b_{j+1} = 2\,b_j .
\end{equation*}

\begin{boitealerte}[Statut exact de cette récurrence]
La formule ci-dessus est \emph{exacte} pour $n = 10, 11, 12, 13$ --- les
quatre valeurs de la table de référence. \textbf{Elle devient fausse
dès $n = 14$~} : la règle canonique donne $23\,936$ là où la valeur
exacte est $23\,808$ (écart $+128$), puis $+896$ en $n = 15$,
$+2\,816$ en $n = 16$, $+6\,784$ en $n = 17$ et $+14\,592$ en $n = 18$.
Elle est donc \emph{une description de la table}, \emph{non} une
loi générale. La fonction qui suit ne s'appuie pas sur elle.
\end{boitealerte}
